\subsection{未来工作}

未来工作应继续围绕机器人基础软件栈展开，并进一步上升到智能泛在操作系统的视角。
本文的实验说明，LeRobot 在线控制并不是单纯的模型推理问题，
而是由 AI 推理、反馈控制、硬件 I/O、用户态运行时和 Linux 调度共同构成的
人机物融合软件系统。因此，后续研究可以参考智能泛在操作系统中
``系统架构--任务调度--驱动生态''的研究主线，
将当前针对单一 LeRobot 平台的性能剖析扩展为面向边缘机器人设备的系统化构造方法。

首先，在系统架构层面，可以围绕机器人实时控制负载探索外核式或类外核式的运行支撑机制。
当前系统仍建立在通用 Linux 宏内核和用户态 Python 运行时之上，
相机、串口、推理服务和控制循环需要通过多层软件抽象间接共享有限 CPU 资源。
后续可以以树莓派 5 这类 ARM 边缘平台为原型，
研究安全直通访问、统一内存视图和低开销用户态资源管理机制，
使相机缓冲区、图像张量、推理输入输出和执行器状态能够在受控权限下高效传递。
在不牺牲安全隔离的前提下，这类机制有望减少不必要的数据复制、上下文切换和运行时抖动，
为机器人在线控制提供比普通应用进程更明确的实时性保障。
同时，外核架构并不意味着简单绕过操作系统，
而是需要把资源授权、隔离和安全检查保留在内核侧，
把与机器人任务强相关的资源使用策略下放到用户态运行时或机器人应用中。

其次，在智能任务与系统调度层面，需要从本文的单模型异步推理扩展到多模型协同推理和反馈控制调度。
后续可建立面向机器人在线控制的智能任务抽象，
描述视觉编码、动作策略、语言指令、安全监测和任务规划等模型之间的数据依赖、
时限约束、动作新鲜度要求和融合策略。
在树莓派 5 这类无 GPU 平台上，重点仍应放在 ARM CPU 推理运行时与 Linux 调度策略的协同优化；
在后续具备 NPU、GPU 或多处理器资源的平台上，
则可进一步研究不同模型在 CPU、NPU 和其他加速器之间的分配、迁移与决策融合。
这一路线的目标不是只提高单次 benchmark 的峰值速度，
而是回答在有限计算资源上推理任务、控制任务和 I/O 任务应如何协同调度，
使系统性能能够服务于反馈控制周期、动作队列深度和任务安全约束。

再次，在硬件 I/O 和驱动生态方面，应从具体设备库优化进一步走向用户态驱动框架与泛在资源统一表征。
本文对摄像头和 SCServo 舵机链路的分析表明，
机器人硬件通信往往同时受软件路径、协议语义和物理总线约束影响。
因此，未来工作仍需区分``可以被软件隐藏的等待''和``由协议物理约束决定的等待''：
摄像头路径可继续从 V4L2 缓冲队列、OpenCV 后台线程、MJPEG 解码和图像张量转换等层面分析端到端时延，
明确后台预取在不同帧率与分辨率下的边界；
舵机路径则应围绕 SCServo 半双工总线、串口波特率、同步读写协议和潜在的流水线读机制继续建模，
判断哪些优化只是移动阻塞位置，哪些优化能够真正减少主控制路径耗时。
在此基础上，可进一步抽象不同机器人设备的操作原语，
形成面向传感器、相机、执行器和通信总线的统一接口与用户态驱动 SDK。
对于能够安全直通的设备，可探索寄存器映射、中断通知和 DMA 缓冲管理等用户态驱动机制；
对于通过 USB、串口等间接接入的设备，则可从协议描述、状态机、同步/异步语义和错误恢复策略出发建立统一表征。
这种驱动框架还可以结合大模型辅助代码生成与迁移工具，
降低新设备接入和现有驱动重构的成本，
使机器人基础软件逐步具备开放、可扩展的硬件生态。

最后，本文使用的跨层 profiling 方法仍可进一步工具化，并与任务级行为质量评价结合。
后续可以将 Python 日志、\texttt{torch.profiler}、\texttt{ftrace}、
内核自定义插桩和实验 CSV 自动对齐为统一时间线，
形成面向机器人基础软件的可复用性能诊断工具。
该工具应能够自动标注每一帧中的观测、推理、执行和等待区间，
关联用户态函数、内核阻塞点和硬件设备事件，
并输出适合优化决策的指标，例如推理临界时间、stall 触发原因、
I/O 等待占比和调度抖动来源。
同时，还需要将这些系统级指标与抓取成功率、末端轨迹误差、动作平滑度、
安全裕量和故障降级次数等任务级指标关联起来，
避免只优化 FPS 而忽略机器人行为质量。
通过这种方式，本文的工作可以从一次针对 LeRobot 的系统剖析，
进一步发展为面向边缘机器人和智能泛在操作系统的通用评测、构造与优化方法。
